\documentclass[11pt,a4paper]{article}

\usepackage[margin=2.3cm]{geometry}
\usepackage{amsmath,amssymb}
\usepackage{booktabs}
\usepackage{xcolor}
\usepackage[most]{tcolorbox}
\usepackage[colorlinks=true,linkcolor=blue!55!black,citecolor=blue!55!black,
            urlcolor=blue!55!black]{hyperref}

\newcommand{\dalpha}{\Delta\alpha}
\newcommand{\ReG}{\operatorname{Re}G_{\perp}}
\newcommand{\ee}{\mathrm{e}}
\newcommand{\ii}{\mathrm{i}}
\newcommand{\dd}{\mathrm{d}}
\newcommand{\unit}[1]{\,\mathrm{#1}}

\title{\bfseries Challenge 20: Torsional levitated optomechanics\\[2pt]
\large Derivation of the librational frequency $\omega_t$ and of the
light--induced torsion--torsion coupling $g$}
\author{}
\date{}

\begin{document}
\maketitle
\vspace{-2.2\baselineskip}

\begin{center}
\small
Companion implementation: \texttt{challenges/20/answer.py}
\end{center}

\section{Statement and strategy}

Two identical dielectric prolate ellipsoids (semi-major axis $a$, two equal
semi-minor axes $b$, relative permittivity $\epsilon_r$, mass density $\rho$) are
levitated in two Gaussian optical tweezers that propagate along $\hat z$, are
linearly polarised along $\hat x$, carry power $P_0$, have waist $w_0$ and wave
vector $k$, and whose foci are separated by $R$ along $\hat x$.  The long axes
align with the polarisation and perform small librations about it.  We must
produce $\omega_t$ and $g$ in

\begin{equation}
  H=\hbar\omega_t\left(a_1^{\dagger}a_1+a_2^{\dagger}a_2\right)
   +\hbar g\left(a_1^{\dagger}a_2+a_1a_2^{\dagger}\right).
  \label{eq:H}
\end{equation}

The derivation has four steps: (i) the quasi-static polarisability tensor of a
sub-wavelength ellipsoid; (ii) the orientational potential of one particle in its
tweezer, giving $\omega_t$; (iii) the light-induced (optical-binding)
dipole--dipole energy expanded to bilinear order in the two libration angles;
(iv) canonical quantisation plus the rotating-wave approximation, giving $g$.

\section{Polarisability of the ellipsoid}

For $a,b\ll\lambda$ the particle responds as a point dipole with the
quasi-static polarisability tensor
$\boldsymbol{\alpha}=\mathrm{diag}(\alpha_{\parallel},\alpha_{\perp},\alpha_{\perp})$
in the body frame \cite{Landau,Osborn,BH},
\begin{equation}
  \alpha_j=\frac{\epsilon_0V(\epsilon_r-1)}{1+L_j(\epsilon_r-1)},
  \qquad V=\frac{4\pi}{3}ab^{2},
  \label{eq:alpha}
\end{equation}
with depolarisation factors, for eccentricity $e=\sqrt{1-b^{2}/a^{2}}$,
\begin{equation}
  L_{\parallel}=\frac{1-e^{2}}{e^{3}}
  \left[\frac{1}{2}\ln\!\frac{1+e}{1-e}-e\right],
  \qquad
  L_{\perp}=\frac{1-L_{\parallel}}{2}.
  \label{eq:L}
\end{equation}
The single quantity that controls both answers is the polarisability anisotropy
\begin{equation}
  \dalpha\equiv\alpha_{\parallel}-\alpha_{\perp}
  =\epsilon_0V(\epsilon_r-1)^{2}
   \frac{L_{\perp}-L_{\parallel}}
        {\bigl[1+L_{\parallel}(\epsilon_r-1)\bigr]
         \bigl[1+L_{\perp}(\epsilon_r-1)\bigr]}>0 .
  \label{eq:dalpha}
\end{equation}

\section{Librational frequency $\omega_t$}

Let $\hat n$ be the long axis and $\theta$ the angle between $\hat n$ and the
polarisation $\hat x$.  Decomposing $\mathbf E=E_0\hat x$ along the principal
axes, the cycle-averaged interaction energy of the induced dipole is
\begin{equation}
  U(\theta)=-\tfrac14 E_0^{2}
  \left[\alpha_{\parallel}\cos^{2}\theta+\alpha_{\perp}\sin^{2}\theta\right]
  =\mathrm{const}+\tfrac14\dalpha E_0^{2}\sin^{2}\theta ,
  \label{eq:Utheta}
\end{equation}
which is Eq.~(1) of Hoang \emph{et al.}~\cite{Hoang2016} written with
$\tilde\chi_j=\alpha_j/(\epsilon_0V)$.  Expanding to $\mathcal O(\theta^{2})$
gives the angular spring constant
\begin{equation}
  k_{\rm tor}=\frac{\partial^{2}U}{\partial\theta^{2}}\bigg|_{0}
  =\frac{1}{2}\dalpha E_0^{2}=\frac{\dalpha I_0}{\epsilon_0c},
  \qquad
  I_0=\frac{2P_0}{\pi w_0^{2}},
  \qquad
  E_0^{2}=\frac{2I_0}{\epsilon_0c}=\frac{4P_0}{\pi w_0^{2}\epsilon_0c}.
  \label{eq:ktor}
\end{equation}
Here $I_0$ is the peak intensity of the Gaussian tweezer and $E_0$ the field
\emph{amplitude} at the focus.  Because the spheroid is symmetric about $\hat n$,
the two librational planes containing $\hat x$ are degenerate and rotation about
$\hat n$ is free \cite{Stickler2021}: there is exactly one torsional mode per
particle, as assumed in Eq.~\eqref{eq:H}.  With mass $m=\rho V$ and the transverse centroidal moment of
inertia
\begin{equation}
  I=\frac{m\left(a^{2}+b^{2}\right)}{5}
   =\frac{4\pi}{15}\rho\,ab^{2}\left(a^{2}+b^{2}\right),
  \label{eq:inertia}
\end{equation}
the librational frequency is $\omega_t=\sqrt{k_{\rm tor}/I}$, i.e.

\begin{tcolorbox}[colback=blue!4,colframe=blue!45!black,title=Result 1:
librational frequency]
\begin{equation}
  \boxed{\;
  \omega_t=\sqrt{\frac{\dalpha\,I_0}{\epsilon_0 c\,I}}
         =\sqrt{\frac{2\,\dalpha\,P_0}{\pi\epsilon_0 c\,w_0^{2}\,I}}
         =\sqrt{\frac{15\,\dalpha\,P_0}
                     {2\pi^{2}\epsilon_0c\,\rho\,w_0^{2}ab^{2}(a^{2}+b^{2})}}\;}
  \label{eq:omegat}
\end{equation}
\end{tcolorbox}

\noindent
Written with the normalised anisotropy
$\Delta\tilde\chi=\dalpha/(\epsilon_0V)$ this is
$\omega_t=\bigl[10\,\Delta\tilde\chi P_0/(\pi\rho c\,w_0^{2}(a^{2}+b^{2}))\bigr]^{1/2}$,
algebraically identical to Eq.~(2) of Ref.~\cite{Hoang2016}, which was verified
experimentally through its $\omega_t\propto\sqrt{P_0}$ scaling.

\section{Light-induced dipole--dipole coupling}

\subsection{Tilt-induced transverse dipoles}

For a tilt $\theta_i\ll1$ of particle $i$ in a plane containing $\hat x$ (unit
vector $\hat e_{\perp}\perp\hat x$ in that plane; $\hat e_\perp=\hat y$ or
$\hat z$ --- both give the same result, see Sec.~\ref{sec:geometry}), the induced
dipole is
\begin{equation}
  \mathbf p_i=\alpha_{\perp}E_0\hat x+\dalpha E_0\cos\theta_i\,\hat n_i
  =E_0\left[\left(\alpha_{\parallel}-\dalpha\theta_i^{2}\right)\hat x
  +\dalpha\,\theta_i\,\hat e_{\perp}\right]+\mathcal O(\theta_i^{3}).
  \label{eq:pi}
\end{equation}

\subsection{Interaction energy}

Two oscillating dipoles separated by $\mathbf R=R\hat x$ interact through the
retarded dipole field.  Writing the instantaneous energy $-\mathbf p_1\cdot
\mathbf E_2(\mathbf r_1)$ with the free-space dyadic Green function
\cite{Dholakia2010,Rieser2022,Rudolph2024},
\begin{equation}
  U_{\rm int}=C\,\operatorname{Re}\!\left\{
  \frac{\ee^{\ii kR}}{4\pi\epsilon_0R^{3}}
  \Bigl[\bigl(\mathbf p_1\!\cdot\!\mathbf p_2
        -3(\mathbf p_1\!\cdot\!\hat R)(\mathbf p_2\!\cdot\!\hat R)\bigr)
        (1-\ii kR)
        -k^{2}R^{2}(\mathbf p_1\!\times\!\hat R)\!\cdot\!
                   (\mathbf p_2\!\times\!\hat R)\Bigr]\right\}.
  \label{eq:Uint}
\end{equation}
The prefactor $C$ is discussed in the boxed remark below: $C=\tfrac12$ for the
correct cycle average, $C=1$ in the convention of Ref.~\cite{Zhang2025}.

Insert Eq.~\eqref{eq:pi} with $\hat R=\hat x$ and collect the terms bilinear in
$\theta_1\theta_2$.  Since the transverse components are perpendicular to
$\hat R$, they do not appear in $(\mathbf p_i\cdot\hat R)$, so the middle term
contributes nothing; the remaining two give
$\mathbf p_1\cdot\mathbf p_2\supset\dalpha^{2}E_0^{2}\theta_1\theta_2$ and
$(\mathbf p_1\times\hat R)\cdot(\mathbf p_2\times\hat R)\supset
\dalpha^{2}E_0^{2}\theta_1\theta_2$.  Hence
\begin{equation}
  U_{\rm int}^{(\theta_1\theta_2)}
  =C\,\dalpha^{2}E_0^{2}\,\theta_1\theta_2\,
   \operatorname{Re}\!\left[\frac{\ee^{\ii kR}\left(1-\ii kR-k^{2}R^{2}\right)}
                                 {4\pi\epsilon_0R^{3}}\right]
  =-C\,\dalpha^{2}E_0^{2}\,\ReG\,\theta_1\theta_2 ,
  \label{eq:Ubilinear}
\end{equation}
where the transverse component of the Green function and its real part are
\begin{equation}
  G_{\perp}=\frac{\ee^{\ii kR}\left(k^{2}R^{2}+\ii kR-1\right)}
                 {4\pi\epsilon_0R^{3}},
  \qquad
  \ReG=\frac{\left(k^{2}R^{2}-1\right)\cos kR-kR\sin kR}{4\pi\epsilon_0R^{3}} .
  \label{eq:G}
\end{equation}
Note that $\ReG$ decays only as $1/R$ in the far field, so the torsion--torsion
interaction is long-ranged and oscillates in sign with $kR$.

\subsection{Quantisation}

The two librators plus Eq.~\eqref{eq:Ubilinear} give
\begin{equation}
  H=\sum_{i=1,2}\left[\frac{L_i^{2}}{2I}
    +\frac{1}{2}k_{\rm tor}\theta_i^{2}\right]+\kappa\,\theta_1\theta_2,
  \qquad
  \kappa=-C\dalpha^{2}E_0^{2}\,\ReG .
\end{equation}
With $\theta_i=\theta_{\rm zpf}(a_i+a_i^{\dagger})$ and
$\theta_{\rm zpf}=\sqrt{\hbar/2I\omega_t}$, and dropping the counter-rotating
$a_1a_2+\mathrm{h.c.}$ terms (valid since $|g|\ll\omega_t$), one obtains
Eq.~\eqref{eq:H} with $\hbar g=\kappa\,\theta_{\rm zpf}^{2}$:

\begin{tcolorbox}[colback=blue!4,colframe=blue!45!black,title=Result 2:
torsion--torsion coupling]
\begin{equation}
  \boxed{\;
  g=-\,C\,\frac{\dalpha^{2}E_0^{2}\,\ReG}{2I\omega_t}
   =-\,C\,\frac{\dalpha^{2}P_0
     \left[\left(k^{2}R^{2}-1\right)\cos kR-kR\sin kR\right]}
     {2\pi^{2}\epsilon_0^{2}c\,w_0^{2}R^{3}\,I\,\omega_t}\;}
  \label{eq:g}
\end{equation}
with $I$ from Eq.~\eqref{eq:inertia}, $\omega_t$ from Eq.~\eqref{eq:omegat},
$E_0^{2}=4P_0/(\pi w_0^{2}\epsilon_0c)$, and $C=\tfrac12$ ($C=1$ in the
convention of Ref.~\cite{Zhang2025}).
\end{tcolorbox}

\noindent
Because $\dalpha^{2}E_0^{2}/(2I)=k_{\rm tor}\dalpha=\omega_t^{2}I\dalpha$,
Eq.~\eqref{eq:g} collapses to the compact, power-independent ratio
\begin{equation}
  \frac{g}{\omega_t}=-\,C\,\dalpha\,\ReG ,
  \label{eq:ratio}
\end{equation}
a useful check: both $g$ and $\omega_t$ scale as $\sqrt{P_0}$, so their ratio is
a pure function of geometry and $kR$.  The normal modes of Eq.~\eqref{eq:H} are
the symmetric/antisymmetric librations at $\omega_{\pm}=\omega_t\pm g$.

\begin{tcolorbox}[colback=red!5,colframe=red!55!black,
  title=\textbf{FLAG --- the factor of two in $g$}]
Equation~\eqref{eq:Uint} with $C=1$ is the \emph{instantaneous} electrostatic-like
energy $U=-\mathbf p_1\cdot\mathbf E_2$ evaluated with field \emph{amplitudes}.
The dipoles oscillate at the optical frequency, so the physical interaction is the
cycle average, which carries an extra factor $\tfrac12$: $C=\tfrac12$.  This is
the convention of the supplementary material of Rieser \emph{et al.}
\cite{Rieser2022}, whose optical-binding force
$\mathbf F_{\rm bind}=\alpha_j\nabla_j\operatorname{Re}
\bigl[(\alpha_{j'}/2\epsilon_0)\,\mathbf E_0^{*}(\mathbf r_j)\cdot
\mathsf G\cdot\mathbf E_0(\mathbf r_{j'})\bigr]$
is exactly the gradient of $-\tfrac{1}{2\epsilon_0}
\operatorname{Re}[\mathbf p_1^{*}\!\cdot\!\mathsf G\!\cdot\!\mathbf p_2]$, and of
the quantum treatment of Rudolph \emph{et al.}~\cite{Rudolph2024,Rudolph2024b}.

Reference~\cite{Zhang2025}, from which this problem is taken, uses $C=1$ in its
Eqs.~(4)--(9) while it \emph{does} cycle-average the trapping term
($k_{\rm tor}=\dalpha E_0^{2}/2$, our Eq.~\eqref{eq:ktor}).  Its $g_0$ is
therefore a factor of two larger than the cycle-averaged result.  Both are
implemented in \texttt{answer.py} via the \texttt{convention} keyword; the
cross-check of Sec.~\ref{sec:crosscheck} is quoted with $C=1$ so that it compares
like with like.  $\omega_t$ is unaffected by this ambiguity.
\end{tcolorbox}

\section{Remarks on the geometry and on neglected terms}
\label{sec:geometry}

\paragraph{Polarisation parallel vs.\ perpendicular to $\mathbf R$.}
The problem places both the polarisation and the separation along $\hat x$
(head-to-tail dipoles), whereas Ref.~\cite{Zhang2025} polarises along $\hat y$
with the separation along $\hat x$ (side-by-side dipoles).  This does \emph{not}
change $g$: in both cases the tilt-induced transverse dipole components are
perpendicular to $\hat R$, so they drop out of $(\mathbf p_i\cdot\hat R)$ and the
bilinear term of Eq.~\eqref{eq:Ubilinear} is the same.  Only the
$\theta_i^{2}$ terms differ between the two geometries.

\paragraph{Binding-induced frequency shift.}
The $\mathcal O(\theta_i^{2})$ part of Eq.~\eqref{eq:Uint} renormalises
$k_{\rm tor}$.  In the head-to-tail geometry of this problem the
$\theta_i^{2}$ terms of $\mathbf p_1\cdot\mathbf p_2-3(\mathbf p_1\cdot\hat R)
(\mathbf p_2\cdot\hat R)$ give a relative correction
$4\alpha_{\parallel}\left[\cos kR+kR\sin kR\right]/(4\pi\epsilon_0R^{3})
\approx2\%$ to $\omega_t^{2}$ (i.e.\ $\sim1\%$ to $\omega_t$) at $R\simeq\lambda$
for the parameters of Table~\ref{tab:numbers}.  It is neglected in
Eq.~\eqref{eq:omegat}, consistent with Ref.~\cite{Zhang2025}.

\paragraph{Beyond the dipole approximation.}
For the benchmark parameters $a/\lambda=0.28$, so quadrupole and higher
multipoles are not negligible; Appendix~A of Ref.~\cite{Zhang2025} estimates
corrections of a few percent.  Likewise, optical binding is genuinely
non-reciprocal and non-Hermitian \cite{Rieser2022,Rudolph2024}; the Hermitian
beam-splitter form assumed in Eq.~\eqref{eq:H} is the conservative part only.

\section{Numerical evaluation}

Table~\ref{tab:numbers} evaluates the results for the parameters of
Ref.~\cite{Zhang2025} ($a=300$~nm, $b=180$~nm, $\epsilon_r=5.7$,
$\rho=3500\unit{kg\,m^{-3}}$, $w_0=500$~nm, $\lambda=1064$~nm, $R=1.06\,\mu$m).

\begin{table}[h]
\centering
\begin{tabular}{lll}
\toprule
Quantity & This work & Ref.~\cite{Zhang2025} \\
\midrule
$V$                       & $4.0715\times10^{-20}\unit{m^{3}}$ & $0.04\times10^{-18}\unit{m^{3}}$ \\
$L_{\parallel}$           & $0.20996$                          & $0.210$ \\
$L_{\perp}$               & $0.39502$                          & $0.395$ \\
$\dalpha$                 & $2.5966\times10^{-31}\unit{C\,m^{2}\,V^{-1}}$ & --- \\
$I$                       & $3.4885\times10^{-30}\unit{kg\,m^{2}}$ & $3.49\times10^{-30}\unit{kg\,m^{2}}$ \\
$\omega_t/2\pi$ (600~mW)  & $1.0418$~MHz & --- \\
$\omega_t/2\pi$ (1000~mW) & $1.3449$~MHz & --- \\
$|g|/2\pi$ (1000~mW, $C=\tfrac12$) & $50.5$~kHz & --- \\
$|g|/2\pi$ (1000~mW, $C=1$)        & $101.0$~kHz & $\approx119$~kHz (quoted) \\
$|g|/\omega_t$ ($C=1$)             & $0.0751$    & --- \\
\bottomrule
\end{tabular}
\caption{Derived quantities for the benchmark parameters.}
\label{tab:numbers}
\end{table}

\section{Cross-check against the published data}
\label{sec:crosscheck}

The authors of Ref.~\cite{Zhang2025} deposited the data behind their Fig.~2(b)
--- $g_0/2\pi$ versus $R/\lambda$ at three tweezer powers --- on Zenodo
\cite{Zenodo}.  Table~\ref{tab:crosscheck} compares those numbers with
Eq.~\eqref{eq:g} evaluated with $C=1$ (their convention).

\begin{table}[h]
\centering
\small
\begin{tabular}{r rr rr rr r}
\toprule
& \multicolumn{2}{c}{600~mW} & \multicolumn{2}{c}{800~mW}
& \multicolumn{2}{c}{1000~mW} & \\
\cmidrule(lr){2-3}\cmidrule(lr){4-5}\cmidrule(lr){6-7}
$R/\lambda$ & Zenodo & Eq.~\eqref{eq:g} & Zenodo & Eq.~\eqref{eq:g}
            & Zenodo & Eq.~\eqref{eq:g} & ratio \\
\midrule
0.550 & $-151.24$ & $+113.29$ & $-174.63$ & $+130.81$ & $-195.24$ & $+146.25$ & $-1.3350$ \\
0.700 & $-11.67$ & $+8.74$ & $-13.48$ & $+10.10$ & $-15.07$ & $+11.29$ & $-1.3350$ \\
0.850 & $+89.93$ & $-67.37$ & $+103.85$ & $-77.79$ & $+116.10$ & $-86.97$ & $-1.3350$ \\
0.950 & $+109.29$ & $-81.87$ & $+126.20$ & $-94.54$ & $+141.10$ & $-105.69$ & $-1.3350$ \\
0.996 & $+104.46$ & $-78.25$ & $+120.63$ & $-90.36$ & $+134.86$ & $-101.02$ & $-1.3350$ \\
1.200 & $+15.73$ & $-11.78$ & $+18.16$ & $-13.60$ & $+20.31$ & $-15.21$ & $-1.3350$ \\
1.500 & $-70.11$ & $+52.52$ & $-80.96$ & $+60.65$ & $-90.52$ & $+67.80$ & $-1.3350$ \\
2.000 & $+52.85$ & $-39.59$ & $+61.02$ & $-45.71$ & $+68.23$ & $-51.11$ & $-1.3350$ \\
3.000 & $+35.36$ & $-26.48$ & $+40.83$ & $-30.58$ & $+45.64$ & $-34.19$ & $-1.3350$ \\
5.000 & $+21.25$ & $-15.92$ & $+24.54$ & $-18.38$ & $+27.44$ & $-20.55$ & $-1.3350$ \\
\bottomrule
\end{tabular}
\caption{$g/2\pi$ in kHz.  ``Zenodo'' is the tabulated Fig.~2(b) data of
Ref.~\cite{Zenodo}; ``Eq.~\eqref{eq:g}'' is this derivation with $C=1$.  The last
column is Zenodo\,/\,ours at 1000~mW.  Generated by \texttt{answer.py}.}
\label{tab:crosscheck}
\end{table}

Three observations:
\begin{enumerate}\itemsep2pt
\item The \textbf{$R$ dependence is reproduced exactly}.  Over
$0.55\le R/\lambda\le6.0$ ($2725$ tabulated points, of which ten are shown) the
ratio is constant to $15$ significant digits, so
Eqs.~\eqref{eq:G}--\eqref{eq:g} carry the correct retarded Green function,
including the sign oscillations and the $1/R$ far-field tail.
\item The \textbf{power dependence is reproduced exactly}: the tabulated values
scale as $\sqrt{P_0}$ (e.g.\ $109.29/141.10=0.7745=\sqrt{0.6}$), confirming that
the $1/\omega_t$ of Eq.~\eqref{eq:g} --- and not $g\propto E_0^{2}\propto P_0$ as
stated in the text of Ref.~\cite{Zhang2025} --- is what their code implements.
\item Two constant discrepancies remain.  The \textbf{overall sign} is a
convention: Ref.~\cite{Zhang2025} writes $g_0=+\dalpha^{2}E_0^{2}\ReG/
(2\sqrt{I_1I_2\omega_1\omega_2})$, whereas Eq.~\eqref{eq:H} with the potential
\eqref{eq:Ubilinear} fixes the opposite sign.  The residual \textbf{factor
$1.334977$} is not accounted for by anything in the published equations: it is
not the multipole correction of their Appendix~A (their own $C_1+C_2$ gives only
$1.045$), not rounding of $L_{\parallel},L_{\perp}$ ($-0.06\%$), and not the
choice of $c$ ($-0.03\%$); $\epsilon_0$ cancels from $g$ identically.  It appears
to be an undocumented constant inside their plotting code.  Their quoted
$119$~kHz lies between our $101$~kHz and their own curve's $135$~kHz at the same
$R$ and $P_0=1$~W.
\end{enumerate}

Independently, $\omega_t$ of Eq.~\eqref{eq:omegat} is algebraically identical to
the experimentally validated expression of Hoang \emph{et al.}~\cite{Hoang2016},
which is a second, independent confirmation of that half of the answer.

\begin{thebibliography}{9}
\bibitem{Zhang2025} G.~Zhang, H.~Zhang and Z.~Yin,
\emph{Scalable universal quantum gates between nitrogen-vacancy centers in
levitated nanodiamond arrays}, Phys.\ Rev.\ A \textbf{112}, 032611 (2025),
\href{https://arxiv.org/abs/2504.08194}{arXiv:2504.08194}.

\bibitem{Zenodo} G.~Zhang, H.~Zhang and Z.~Yin, \emph{Quantum gates between
nitrogen-vacancy centers in levitated nanodiamond arrays} (dataset),
\href{https://doi.org/10.5281/zenodo.16917848}{10.5281/zenodo.16917848} (2025).

\bibitem{Hoang2016} T.~M.~Hoang, Y.~Ma, J.~Ahn, J.~Bang, F.~Robicheaux,
Z.-Q.~Yin and T.~Li, \emph{Torsional optomechanics of a levitated nonspherical
nanoparticle}, Phys.\ Rev.\ Lett.\ \textbf{117}, 123604 (2016).

\bibitem{Rieser2022} J.~Rieser \emph{et al.}, \emph{Tunable light-induced
dipole-dipole interaction between optically levitated nanoparticles},
Science \textbf{377}, 987 (2022).

\bibitem{Rudolph2024} H.~Rudolph, U.~Deli\'c, K.~Hornberger and B.~A.~Stickler,
\emph{Quantum theory of non-Hermitian optical binding between nanoparticles},
Phys.\ Rev.\ A \textbf{110}, 063507 (2024).

\bibitem{Rudolph2024b} H.~Rudolph, U.~Deli\'c, K.~Hornberger and
B.~A.~Stickler, \emph{Quantum optical binding of nanoscale particles},
Phys.\ Rev.\ Lett.\ \textbf{133}, 233603 (2024).

\bibitem{Dholakia2010} K.~Dholakia and P.~Zem\'anek, \emph{Colloquium: Gripped by
light --- optical binding}, Rev.\ Mod.\ Phys.\ \textbf{82}, 1767 (2010).

\bibitem{Landau} L.~D.~Landau and E.~M.~Lifshitz,
\emph{Electrodynamics of Continuous Media}, 2nd ed., Sec.~4.

\bibitem{Osborn} J.~A.~Osborn, \emph{Demagnetizing factors of the general
ellipsoid}, Phys.\ Rev.\ \textbf{67}, 351 (1945).

\bibitem{BH} C.~F.~Bohren and D.~R.~Huffman, \emph{Absorption and Scattering of
Light by Small Particles}, Ch.~5 (Wiley, 1983).

\bibitem{Stickler2021} B.~A.~Stickler, K.~Hornberger and M.~S.~Kim,
\emph{Quantum rotations of nanoparticles}, Nat.\ Rev.\ Phys.\ \textbf{3}, 589
(2021).
\end{thebibliography}

\end{document}
